\section{Source audit, integration decisions, and further research}
\label{audit:integration-section}

\subsection{The inspected baseline and the scope of the audit}
\label{audit:baseline}

The comparison baseline is the repository revision
\[
 \texttt{446dca9bd65c20e409d1d2561d4ffad408103e2a}.
\]
The two locations named in the research request are the canonical
polylogarithm manuscript and the current incoming directory
\cite{Repo,Incoming}. The six current incoming archives were extracted
and their relevant statements compared with the manuscript's overview
and proof-status records. The targeted mathematical checks concerned
the even-tail coefficient conventions and operator reduction, the
differential-polynomial normal form and known Gamma coefficient sectors,
the normalized residue calculus, and the first determinantal resonance
and its stated continuation problem \cite{DPL,PTI,BTD,DHR}.

The new arguments were checked against those inputs, including the
omitted binomial summand, the signs and index shifts in the boundary
determinant, the Euler-operator projection, and the endpoint constants
in the convergent integral formulas. This is a targeted audit of the
proof chain used here. It is not a certification of every chapter,
historical archive, or earlier numerical calculation in the repository.
\textbf{No new false inherited theorem was established by this audit.}
The source changes below concern proved extensions, attribution, and
precise status updates rather than a newly discovered counterexample.

\subsection{What is new relative to that baseline}
\label{audit:status-matrix}

The following distinctions prevent an already available reduction from
being counted as a new resolution. ``New'' here means established in
this continuation relative to the inspected source material; it does
not assert historical priority over the entire literature.

\begin{center}
\small
\begin{tabular}{@{}>{\raggedright\arraybackslash}p{0.18\linewidth}>{\raggedright\arraybackslash}p{0.34\linewidth}>{\raggedright\arraybackslash}p{0.40\linewidth}@{}}
\hline
Topic & Inherited input & Exact advance in this article\\
\hline
Even-tail rank & Equal and adjacent orders; the $p=2,4,6$ families;
the fixed-gap boundary determinant in DPL,
\texttt{ar:gap-det} \cite{DPL,PTI,BTD}.
& Strict separation of every pair of recurrence-column ratios, and
nonvanishing of the complete $q-p=4$ family. The fixed-gap determinant
is rederived as an input, not claimed as new.\\[4pt]
Finite telescopers & The Bernoulli coordinate formula and explicit
low-order telescopers, including $(4,6)$ \cite{DPL}.
& A polynomial certificate for every late-block combination and every
even pair, with a uniform construction, degree bounds, and limiting
constant. The $(8,12)$ rational extraction is an explicit application.\\[4pt]
Gamma coefficients & The canonical quotient, Appell transport, the
complete $u_0,u_1$ sector, and every pure polygamma-square coefficient
\cite{DPL,PTI,BTD}.
& The complete projection onto $u_0,u_1,u_2$, including all powers
$u_2^b$, $b\ge2$, and the rational coefficient form with numerator
degree $b-2$. These are formal coefficient identities.\\[4pt]
Next resonance & Classical Gauss summation and the determinant
identity; the first-resonance theory and proposed next two-particle
target \cite{DHR}.
& The next double pole, including overlap subtraction, the hard-cutoff
constant, its spectral comparison, uniform cutoff jets, and anchored
Hurwitz--Stieltjes identities.\\[4pt]
Higher spectral jets & The normalized principal-part expansion,
scalar projection, and first spectral-coefficient bridge \cite{DPL}.
& Canonical spatial kernels for every $[s^q]D_1(s)$, $q\ge1$, with
all origin, infinity, and higher-pole terms. The same scalar projection
controls every Stieltjes correction.\\[4pt]
Family--2 example & An exact incoming proof of the manuscript's
\texttt{double:eq:fam2-example}, with a simplified value \cite{PTI}.
& An editorial promotion supported by that inherited proof. No new
proof of this example is claimed here.\\
\hline
\end{tabular}
\end{center}

\subsection{Concrete integration decisions and limits}
\label{audit:decisions}

The optimal-prefix conjecture should now record the union of the
proved equal-order, adjacent-order, $p=2,4,6$, and gap-four families.
Its general unequal-order statement remains open. In particular,
pairwise ratio separation does not settle a boundary determinant of
size three or larger. There is a simple algebraic reason: the positive
matrix $B_{ij}=i+j$, $1\le i,j\le3$, has every increasing
two-by-two minor strictly negative,
\[
 \det B_{\{i_1,i_2\},\{j_1,j_2\}}
       =-(i_2-i_1)(j_2-j_1)<0,
\]
but its three-by-three determinant is zero. Thus the new ratio theorem
supplies exactly the sign information needed for a two-column
boundary block; an additional argument is required for $m\ge3$.

The rational telescoper theorem changes the status of the requested
all-order certificate construction. It does not by itself prove that
an arbitrary prescribed inverse row exists. When a coefficient relation
has been supplied, the theorem constructs and verifies its finite
certificate; the separate rank theorem proves existence of every
reconstruction row in the stated families.

For the Gamma sector, retain the order of operations: first take the
canonical normal form and then apply the projection $\sigma_2$ which
sets $u_r=0$ for $r\ge3$. The Euler-operator proof justifies extracting
this projected sector; setting higher derivatives to zero before
normal reduction is a different operation. The all-$b$ recurrence and
rational numerator theorem extend the earlier sector results. They
neither evaluate every associated finite-part moment nor establish
independence of the resulting numerical constants \cite{DPL,PTI}.

All normalizations should travel with their formulas. The tail
identities use the centered coordinate $n+1/2$ and explicit polynomial
subtraction. Differential-polynomial finite parts retain the original
spatial coordinate and the endpoint functional of total derivatives.
The logarithmic kernels use the displayed subtraction at infinity and
the splitting point $1$. The resonant hard cutoff uses $x=N+1$, while
its spectral constant belongs to the explicitly stated Gamma
deformation. The product of a vanishing tail with a divergent moment
contributes to that constant and must be retained. These distinctions
are requirements of the proved identities, not interchangeable
definitions of one unspecified regularization.

The canonical Family--2 display
\texttt{chapters/04-depth.tex}, \texttt{double:eq:fam2-example},
can be marked proved by citing the existing incoming component
evaluations \texttt{bip:eq:12}--\texttt{bip:eq:21} and simplification
\texttt{bip:eq:family2-short} \cite{PTI}. This promotion was already
justified before the present work. The Gaussian $S_6$ and the
\emph{current} $S_8$ candidate remain conjectural. The older rejected
$S_8$ vector is a different statement and must not replace the current
candidate in either a proof attempt or a status table \cite{Repo,PTI}.
None of the new rank, Gamma-sector, or resonance results proves those
period identities.

Exact finite algebra and numerical diagnostics serve different
purposes. The rational replays test coefficient transformations and
finite certificates. The separately subtracted quadratures and
complex-parameter comparisons test analytic formulas and endpoint
constants. The infinite-family results rest on the proofs, rather than
on the extent of a computed table; the floating-point diagnostics are
not interval certificates. No proof-assistant formalization is claimed.

\subsection{Twelve concrete questions for the next continuation}
\label{audit:future}

The following targets emphasize exact identities, structural reductions,
and normalized primitives. Each calls for a proof or a specified finite
certificate; a failed numerical relation search would not settle an
independence or nonmembership claim.

\begin{enumerate}
\item \textbf{Settle the next boundary determinant.}
For $m=3$, prove or disprove nonvanishing of
$\det(C_{b+i,j}^{[3]})_{1\le i,j\le3}$ for every $b>3$.
This is the next gap $q-p=6$. The pairwise signs are now known, but
they do not control the three-column determinant. A discrete Wronskian
identity, an exact recurrence for the determinant, or a nonvanishing
arithmetic normalization would supply the missing step.

\item \textbf{Explain the gap-four determinant sequence arithmetically.}
Find a generating function, finite recurrence, or exact denominator
formula for $\Delta_b=u_{b+2}v_{b+1}-u_{b+1}v_{b+2}$.
The recurrence already proves positivity; the new target is a closed
identity or an integer normalization which explains the numerators
of the optimal-prefix determinants and supports exact reconstruction.

\item \textbf{Describe all reconstruction rows without generic inversion.}
For the proved gap-four family, express each inverse row through
finite binomial or Bernoulli sums and the two fundamental recurrence
solutions. Combine it with the universal telescoper to obtain
certificates for prescribed even-zeta coordinates. Determine minimal
weight-polynomial degrees within the specified formal coefficient
system, keeping this separate from arithmetic independence.

\item \textbf{Extend polynomial telescoping to three or more tails.}
Construct a finite difference certificate involving products of every
required proper subset of the tails. Identify precisely which parity
conditions leave only even reciprocal powers and which additional
Euler-sum coordinates survive. The target is an all-order certificate
and its endpoint constant, rather than an unproved extrapolation of
the two-tail cancellation.

\item \textbf{Determine the complete next Gamma projection.}
Apply the Euler-operator method to $\sigma_3$, retaining
$u_0,u_1,u_2,u_3$. Derive a closed coefficient equation with the
canonical support restrictions, then solve the first coupled families
containing $u_3^2$ and arbitrary powers of $u_1,u_2$.
Prove that discarded higher canonical monomials cannot contribute
before extracting this sector.

\item \textbf{Find a closed formula for the all-$b$ sector numerators.}
The rational form, denominator, degree, and leading coefficient of
the polynomials $p_b(k)$ are now established. Seek a finite formula
uniform in both $b$ and $k$, a bivariate generating function, or a
hypergeometric representation. Such a formula should explain exact
vanishing coefficients and coefficient patterns, rather than merely
reproduce the already proved recurrence.

\item \textbf{Evaluate the numerical image of the new Gamma sector.}
Map the moments containing $(\psi')^{2k}(\psi'')^b$ and their
Gamma--Appell transports into convergent free-order multiple-zeta
derivatives. Determine the relations among these images within an
explicitly stated functional-relation family. This would connect the
new formal coefficient theorem with the remaining basis-reduction questions for
higher mixed moments in the incoming reports \cite{DPL,PTI}.

\item \textbf{Introduce rational shifts into the logarithmic kernels.}
Replace the unshifted pole lattice by a fixed Hurwitz lattice and
derive the exact analogue of the scalar correction for
$\gamma_q(a)$. At rational $a$, seek finite formulas in a declared
multiple-Hurwitz or cyclotomic basis. Retain the origin subtraction,
the choice of logarithm, and the transport of the normalization
constant under changes of base point.

\item \textbf{Construct parameter primitives for the higher jets.}
Starting with the normally convergent kernels, integrate in a shift
or Gamma parameter and anchor the primitive at a specified point.
Find recurrences which isolate one prescribed higher Stieltjes
coefficient or one marked spectral derivative. A complete answer
should prove the endpoint limit and identify any Barnes-type
functions needed for the resulting primitive family.

\item \textbf{Treat higher-rank and deeper determinantal resonances.}
At the next general-rank resonance $b-a-c=2r-3$, construct the
subtraction for the divergent two-by-two moment block together with
the finite lower block. Then study the deeper loci
$b-a-c=2r-2-M$. Determine all overlapping escape terms before
comparing hard-cutoff constants with spectral Laurent coefficients,
and prove uniformity sufficient for every fixed parameter jet.

\item \textbf{Find a depth description compatible with finite differences.}
For rational Gamma weights at other positive integer parameters,
derive complete polylogarithm generators of the largest-index
increments and their anchored Euler primitives. Determine which
cutoff harmonic products cancel under finite differencing and when
the endpoint sums reduce to ordinary zeta products. The two-particle
specialization here supplies a test case for a depth invariant that
respects this operation.

\item \textbf{Return to the frozen Gaussian period targets.}
For the exact $S_6$ vector and current $S_8$ candidate, seek a
parameter-dependent polylogarithm or iterated-integral identity whose
proved boundary value is the required relation. Preserve the branch
data and all lower-weight constants. A projection outside a selected
formal relation space is not a disproof of the numerical period
identity, and a small residual is not its proof \cite{Repo}.
\end{enumerate}
